\appendix{Представление графического материала}

Графический материал, выполненный на отдельных листах,
изображен на рисунках А.1–А.10.
\setcounter{числоПлакатов}{0}

\renewcommand{\thefigure}{А.\arabic{figure}} % шаблон номера для плакатов

\begin{landscape}

\begin{плакат}
	\includegraphics[width=0.82\linewidth]{pos1.eps}
	\заголовок{Сведения о ВКРБ}
	\label{pl1:image}
\end{плакат}

\begin{плакат}
	
\includegraphics[width=0.82\linewidth]{pos2.eps}
	\заголовок{Цель и задачи разработки}
	\label{pl2:image}
\end{плакат}

\begin{плакат}
	\includegraphics[width=0.82\linewidth]{pos3.eps}
	\заголовок{Диаграмма прецедентов для мастера}
	\label{pl3:image}
\end{плакат}

\begin{плакат}
	\includegraphics[width=0.82\linewidth]{pos4.eps}
	\заголовок{Диаграмма прецедентов для клиента}
	\label{pl4:image}
\end{плакат}

\begin{плакат}
	\includegraphics[width=0.82\linewidth]{pos5.eps}
	\заголовок{Диаграмма прецедентов для администратора}
	\label{pl5:image}
\end{плакат}

\begin{плакат}
	\includegraphics[width=0.82\linewidth]{pos6.eps}
	\заголовок{Общая архитектура системы}
	\label{pl6:image}
\end{плакат}

\begin{плакат}
	\includegraphics[width=0.82\linewidth]{pos7.eps}
	\заголовок{Модель данных}
	\label{pl7:image}
\end{плакат}

\begin{плакат}
	\includegraphics[width=0.82\linewidth]{pos8.eps}
	\заголовок{Сценарий бронирования}
	\label{pl8:image}
\end{плакат}

\begin{плакат}
	\includegraphics[width=0.82\linewidth]{pos9.eps}
	\заголовок{Реляционная модель базы данных}
	\label{pl9:image}
\end{плакат}

\begin{плакат}
	\includegraphics[width=0.82\linewidth]{pos10.eps}
	\заголовок{Заключение}
	\label{pl10:image}
\end{плакат}

\end{landscape}
